\section{Evaluation}
\label{sec:evaluation}

We evaluate \greyboxfuzz with three goals. First, we measure whether the input formatter can generate structured seeds at a practical speed. Second, we evaluate whether the static taint tracking component can correctly identify cryptographic branches that should be prioritized by the coverage monitor. Third, we evaluate whether \greyboxfuzz can detect cryptographic logic errors in real-world \zksnark applications more efficiently than a baseline differential fuzzing approach.

All experiments were conducted on an Intel Core i7-8700 CPU @ 3.20 GHz with 32 GB of RAM, running 64-bit Ubuntu 22.04 LTS. Unless otherwise stated, each experiment was repeated under the same configuration, and we report the average result.

\subsection{Seed Generation}

The first experiment evaluates the seed-generation efficiency of \greyboxfuzz. Since \greyboxfuzz uses an input formatter to construct seeds that follow the expected input structure of \zksnark programs, this experiment examines whether the formatting process introduces significant runtime overhead. In \greyboxfuzz, the seed structure is derived from the Charm statement provided by the user. Therefore, this experiment evaluates whether Charm-based input information can be used to generate structured seeds at a practical speed.

We compare \greyboxfuzz with \afl on three \zksnark applications: Circom Sudoku \cite{circom_sudoku}, The Age of Navigation \cite{age_of_navigation}, and dapp-zkatm \cite{dapp_zkatm}. For each tool, we measure the total number of generated seeds and the number of malformed seeds under two time budgets, 30 seconds and 60 seconds. A seed is considered malformed if it does not satisfy the expected input format of the target application. This experiment focuses on the cost and quality of seed construction before analyzing the deeper program logic explored by the fuzzer.

\autoref{tab:seed_generation} summarizes the results. Across all three applications, \greyboxfuzz generates slightly fewer seeds than \afl, but the difference is moderate. For example, within 30 seconds, \afl generates 65,906 seeds for Circom Sudoku, while \greyboxfuzz generates 53,827 seeds. Similarly, for The Age of Navigation and dapp-zkatm, \greyboxfuzz generates 55,167 and 55,251 seeds, compared with 68,704 and 68,159 seeds generated by \afl. The same trend is observed under the 60-second budget. These results show that the Charm-based formatter introduces some additional overhead, but the generation speed remains practical.

More importantly, \greyboxfuzz substantially reduces the malformed seed rate. For \afl, the malformed rate is extremely high across all applications, ranging from 94.04\% to 99.98\%. This is because \afl mutates raw bytes without knowledge of the structured input constraints required by \zksnark applications. In contrast, \greyboxfuzz keeps the malformed rate below 5\% in all settings. For example, under the 60-second budget, the malformed rates of \greyboxfuzz are 3.94\% for Circom Sudoku, 3.91\% for The Age of Navigation, and 4.62\% for dapp-zkatm. This indicates that \greyboxfuzz can preserve the expected input structure much more effectively than \afl.

Overall, the results show that \greyboxfuzz trades a small amount of seed-generation speed for a large improvement in seed validity. This tradeoff is beneficial for \zksnark application fuzzing because malformed seeds are more likely to trigger shallow parsing or format-checking failures, while well-formed seeds allow the fuzzer to reach deeper logic related to \zksnark verification and application-level constraints. Therefore, even though \greyboxfuzz generates seeds slightly more slowly than \afl, it provides better effective fuzzing performance by focusing more executions on structurally valid inputs and reducing time spent on trivial malformed-input failures.


\begin{table}[!htp]
\centering
\makebox[\textwidth][c]{%
\rotatebox{90}{
\begin{adjustbox}{max width=0.9\textheight, max totalheight=0.9\textwidth, keepaspectratio}
\begin{minipage}{\textheight}
\centering
\footnotesize
\renewcommand{\arraystretch}{1.75}
\begin{tabular}{cc|ccc|ccc}
\hline
\multicolumn{2}{c|}{\multirow{2}{*}{Execution Time}} & \multicolumn{3}{c|}{30 Seconds} & \multicolumn{3}{c}{60 Seconds} \\ \cline{3-8} 
\multicolumn{2}{c|}{} & \multicolumn{1}{c|}{\begin{tabular}[c]{@{}c@{}}Number of\\ Total Seeds\end{tabular}} & \multicolumn{1}{c|}{\begin{tabular}[c]{@{}c@{}}Number of\\ Malformed Seed\end{tabular}} & Malformed Rate & \multicolumn{1}{c|}{\begin{tabular}[c]{@{}c@{}}Number of\\ Total Seeds\end{tabular}} & \multicolumn{1}{c|}{\begin{tabular}[c]{@{}c@{}}Number of\\ Malformed Seed\end{tabular}} & Malformed Rate \\ \hline
\multicolumn{1}{c|}{\multirow{3}{*}{AFL}} & Circom Sudoku & \multicolumn{1}{c|}{65906} & \multicolumn{1}{c|}{65888} & 99.97\% & \multicolumn{1}{c|}{112729} & \multicolumn{1}{c|}{112711} & 99.98\% \\ \cline{2-8} 
\multicolumn{1}{c|}{} & The Age of Navigation & \multicolumn{1}{c|}{68704} & \multicolumn{1}{c|}{66801} & 97.23\% & \multicolumn{1}{c|}{134908} & \multicolumn{1}{c|}{131114} & 97.19\% \\ \cline{2-8} 
\multicolumn{1}{c|}{} & dapp-zkatm & \multicolumn{1}{c|}{68159} & \multicolumn{1}{c|}{64102} & 94.04\% & \multicolumn{1}{c|}{134492} & \multicolumn{1}{c|}{127209} & 94.58\% \\ \hline
\multicolumn{1}{c|}{\multirow{3}{*}{\greyboxfuzz}} & Circom Sudoku & \multicolumn{1}{c|}{53827} & \multicolumn{1}{c|}{2195} & 4.08\% & \multicolumn{1}{c|}{101983} & \multicolumn{1}{c|}{4018} & 3.94\% \\ \cline{2-8} 
\multicolumn{1}{c|}{} & The Age of Navigation & \multicolumn{1}{c|}{55167} & \multicolumn{1}{c|}{2147} & 3.89\% & \multicolumn{1}{c|}{106727} & \multicolumn{1}{c|}{4175} & 3.91\% \\ \cline{2-8} 
\multicolumn{1}{c|}{} & dapp-zkatm & \multicolumn{1}{c|}{55251} & \multicolumn{1}{c|}{2692} & 4.87\% & \multicolumn{1}{c|}{104940} & \multicolumn{1}{c|}{4851} & 4.62\% \\ \hline
\end{tabular}

\caption{Seed- generation speed and malformed rate comparison among \afl and \greyboxfuzz under different time budgets.}
\label{tab:seed_generation}
\end{minipage}
\end{adjustbox}%
}%
}
\end{table}




% \begin{table*}[t]
% \centering
% \begin{tabular}{cc|ccc|ccc}
% \hline
% \multicolumn{2}{c|}{\multirow{2}{*}{Execution Time}} & \multicolumn{3}{c|}{30 Seconds} & \multicolumn{3}{c}{60 Seconds} \\ \cline{3-8} 
% \multicolumn{2}{c|}{} & \multicolumn{1}{c|}{\begin{tabular}[c]{@{}c@{}}Number of\\ Total Seeds\end{tabular}} & \multicolumn{1}{c|}{\begin{tabular}[c]{@{}c@{}}Number of\\ Malformed Seed\end{tabular}} & Malformed Rate & \multicolumn{1}{c|}{\begin{tabular}[c]{@{}c@{}}Number of\\ Total Seeds\end{tabular}} & \multicolumn{1}{c|}{\begin{tabular}[c]{@{}c@{}}Number of\\ Malformed Seed\end{tabular}} & Malformed Rate \\ \hline
% \multicolumn{1}{c|}{\multirow{3}{*}{AFL}} & Circom Sudoku & \multicolumn{1}{c|}{65906} & \multicolumn{1}{c|}{65888} & 99.97\% & \multicolumn{1}{c|}{112729} & \multicolumn{1}{c|}{112711} & 99.98\% \\ \cline{2-8} 
% \multicolumn{1}{c|}{} & The Age of Navigation & \multicolumn{1}{c|}{68704} & \multicolumn{1}{c|}{66801} & 97.23\% & \multicolumn{1}{c|}{134908} & \multicolumn{1}{c|}{131114} & 97.19\% \\ \cline{2-8} 
% \multicolumn{1}{c|}{} & dapp-zkatm & \multicolumn{1}{c|}{68159} & \multicolumn{1}{c|}{64102} & 94.04\% & \multicolumn{1}{c|}{134492} & \multicolumn{1}{c|}{127209} & 94.58\% \\ \hline
% \multicolumn{1}{c|}{\multirow{3}{*}{\greyboxfuzz}} & Circom Sudoku & \multicolumn{1}{c|}{53827} & \multicolumn{1}{c|}{2195} & 4.08\% & \multicolumn{1}{c|}{101983} & \multicolumn{1}{c|}{4018} & 3.94\% \\ \cline{2-8} 
% \multicolumn{1}{c|}{} & The Age of Navigation & \multicolumn{1}{c|}{55167} & \multicolumn{1}{c|}{2147} & 3.89\% & \multicolumn{1}{c|}{106727} & \multicolumn{1}{c|}{4175} & 3.91\% \\ \cline{2-8} 
% \multicolumn{1}{c|}{} & dapp-zkatm & \multicolumn{1}{c|}{55251} & \multicolumn{1}{c|}{2692} & 4.87\% & \multicolumn{1}{c|}{104940} & \multicolumn{1}{c|}{4851} & 4.62\% \\ \hline
% \end{tabular}
% \caption{Seed- generation speed and malformed rate comparison among \afl and \greyboxfuzz under different time budgets.}
% \label{tab:seed_generation}
% \end{table*}



\subsection{Coverage Evaluation}

The second experiment evaluates the effectiveness of the coverage monitor. As discussed in Section~\ref{sec:coverage_monitor}, \greyboxfuzz uses static taint tracking to identify cryptographic branches before fuzzing begins. The purpose of this experiment is to determine whether the taint tracker can correctly identify branches related to proof generation and proof verification.

We use \rapidsnark as the target library for this evaluation. \rapidsnark is a C++ implementation of \zksnark proof generation and verification for circuits produced by Circom and snarkjs. It is a suitable target because it contains both prover and verifier components and is used in practical \zksnark workflows. To construct a reliable ground truth, we compiled \rapidsnark version v0.0.8 with debugging information enabled. This allows us to correlate assembly instructions with source-code locations and manually determine whether each branch belongs to cryptographic computation or general program logic.

After compilation, we extracted the assembly code of the prover and verifier and manually labeled each branch as either a cryptographic branch or a non-cryptographic branch. A branch is labeled as cryptographic if it directly participates in proof generation, proof verification, witness processing, finite-field arithmetic, elliptic-curve computation, or pairing-related computation. Branches related only to parsing, file handling, memory management, logging, or other general program functionality are labeled as non-cryptographic. This manual labeling serves as the ground truth for evaluating the static taint tracking result.

We then ran the static taint tracker on the same \rapidsnark binaries. The taint tracker automatically marks branches that are influenced by cryptographic inputs, such as the circuit, witness, proving key, proof, public input, and verification key. We compare the automatically tainted branches with the manually labeled ground truth. \autoref{tab:coverage_performance} reports the number of branches in each component, the number of branches marked by static taint tracking, and the correctness of the classification.

The main goal of this evaluation is to measure whether \greyboxfuzz can identify cryptographic branches with high coverage. Missing cryptographic branches is more harmful than marking a small number of non-cryptographic branches, because missed cryptographic branches may prevent the fuzzer from exploring security-relevant computation. In contrast, a small number of false positives only causes \greyboxfuzz to monitor additional branches. Therefore, this evaluation focuses on whether the taint tracker successfully captures the cryptographic branches needed for guiding fuzzing.


The results show that static taint tracking can effectively identify the cryptographic branches needed by the coverage monitor. Only 14 out of 250 \cryptobranch branches (5.6\%) are not successfully identified by the taint tracker, which means that \greyboxfuzz captures most branches related to proof generation and proof verification. Meanwhile, 36 out of 239 \programbranch branches (15.1\%) are incorrectly identified as \cryptobranch branches. However, these false positives do not cause \greyboxfuzz to miss critical vulnerabilities. Instead, they only cause \greyboxfuzz to monitor a small number of additional branches. Since most cryptographic branches are correctly tainted and only a limited number of non-cryptographic branches are misclassified, \greyboxfuzz can guide fuzzing toward proof-related computation without introducing significant unnecessary overhead.



\begin{table*}[!t]
\centering
\footnotesize
\renewcommand{\arraystretch}{1.15}
\begin{tabular}{c|c|ccc|cc}
\toprule
Branch Type & Total Branch & Tainted & Untainted & Mis-Tainted & Error Rate & Correct Rate \\
\midrule
Cryptographic & 250 & 236 & 14 & 14 & 5.6\% & 94.4\% \\
\midrule
Non-cryptographic & 239 & 36 & 203 & 36 & 15.1\% & 84.9\% \\
\midrule
Overall & 489 & 272 & 217 & 50 & 10.2\% & 89.8\% \\
\bottomrule
\end{tabular}
\caption{Coverage evaluation of static taint tracking on \rapidsnark.}
\label{tab:coverage_performance}
\end{table*}


\subsection{Performance Evaluation}

After demonstrating the seed-generation speed and coverage performance of \greyboxfuzz, we run a third experiment to evaluate its bug detection performance. We selected a set of \zksnark applications from both real-world projects and known vulnerability patterns. The selected programs include common circuit-level mistakes, application-level input validation bugs, and implementation errors that may cause a verifier to accept an invalid proof or reject a valid proof.

We selected \zksnark applications from different sources, including prototype demos, academic projects, and live applications. These applications are implemented in different programming languages, including C, Rust, and Go, and are built on different \zksnark frameworks and tools, including \rapidsnark, \bellman, \arkworks, \circom, \gnark, and \snarkjs. This diversity allows us to evaluate \greyboxfuzz under a wide range of implementation settings and demonstrates that \greyboxfuzz is not limited to a specific programming language, framework, or development workflow.


\autoref{tab:bug_evaluation} presents the bugs identified in the tested applications. Because \greyboxfuzz includes an error locator, the table also reports the component that caused the detected bug. This additional information helps developers determine whether the issue comes from the circuit, prover-side conversion, proof generation, verifier-side conversion, or verification.

\begin{table*}[!t]
\centering
\small
\renewcommand{\arraystretch}{1.25}
\begin{tabular}{p{0.37\textwidth}|p{0.13\textwidth}|p{0.11\textwidth}|p{0.14\textwidth}|p{0.11\textwidth}}
\hline
\textbf{Error/Vulnerability in \zksnark Program} 
& \textbf{Library} 
& \textbf{Protocol} 
& \textbf{Location} 
& \textbf{Bug Type} \\ \hline\hline

\multicolumn{5}{c}{\textbf{Potentially Locatable in Current Libraries}} \\ \hline

Multiple Gadget Conflict Issue 
& All libraries 
& R1CS 
& Circuit 
& Common \\ \hline

Incorrect Bit/Comparison Gadget Implementation 
& \libsnark 
& R1CS 
& Circuit 
& Common \\ \hline

Developer Incorrect Flattening 
& All libraries 
& R1CS 
& Circuit 
& Common \\ \hline

BigInt: Missing Bit Length Check
& All libraries 
& R1CS 
& Circuit 
& Common \\ \hline

Arithmetic Over/Under Flows 
& Circom 
& R1CS 
& Circuit 
& Common \\ \hline\hline

\multicolumn{5}{c}{\textbf{Bugs in the Wild}} \\ \hline

Unused Variables Issues in Circuit 
& gnark 
& Groth16 
& Setup/Prover 
& Real-world \\ \hline

\zksnark in Federated Learning
& Circom 
& Groth16 
& Circuit 
& Real-world \\ \hline

\end{tabular}
\caption{Error and vulnerability cases used to evaluate the error localization capability of \greyboxfuzz.}
\label{tab:bug_evaluation}
\end{table*}

The expected benefit of \greyboxfuzz comes from three design choices. First, the input formatter generates structured seeds that are more likely to reach proof generation and verification. Second, the coverage monitor prioritizes seeds that exercise new cryptographic branches instead of treating all execution paths equally. Third, the error locator provides component-level feedback after a mismatch is detected. Together, these features allow \greyboxfuzz to detect cryptographic logic errors more efficiently than a general differential fuzzing baseline.


In Section~\ref{sec:discussion}, we further discuss selected security bugs discovered during the evaluation and explain how the error locator helps identify their causes.